
\documentclass{article}
\usepackage{colortbl}
\usepackage{makecell}
\usepackage{multirow}
\usepackage{supertabular}

\begin{document}

\newcounter{utterance}

\centering \large Interaction Transcript for game `codenames', experiment `ambiguous', episode 9 with reflexion{-}llama{-}3.1{-}8b{-}instruct.
\vspace{24pt}

{ \footnotesize  \setcounter{utterance}{1}
\setlength{\tabcolsep}{0pt}
\begin{supertabular}{c@{$\;$}|p{.15\linewidth}@{}p{.15\linewidth}p{.15\linewidth}p{.15\linewidth}p{.15\linewidth}p{.15\linewidth}}
    \# & $\;$A & \multicolumn{4}{c}{Game Master} & $\;\:$B\\
    \hline

    \theutterance \stepcounter{utterance}  
    & & & \multicolumn{2}{p{0.3\linewidth}}{
        \cellcolor[rgb]{0.9,0.9,0.9}{
            \makecell[{{p{\linewidth}}}]{
                \texttt{\tiny{[GM$|$GM]}}
                {'hidden': {'team': ['angle', 'flush'], 'innocent': ['cuff', 'boil', 'jam', 'fan', 'rocket'], 'opponent': [], 'assassin': ['beetle']}, 'revealed': {'team': {'team': [], 'innocent': [], 'opponent': [], 'assassin': []}, 'opponent': {'team': [], 'innocent': [], 'opponent': [], 'assassin': []}}}
            }
        }
    }
    & & \\ \\

    \theutterance \stepcounter{utterance}  
    & & \multicolumn{4}{p{0.6\linewidth}}{
        \cellcolor[rgb]{0.9,0.9,0.9}{
            \makecell[{{p{\linewidth}}}]{
                \texttt{\tiny{[P1$\langle$GM]}}
                \texttt{You have attempted this same task before and did not succeed. These are the plans you wrote after those attempts.} \\
\texttt{Trial \#1: New plan: In my last attempt, I provided a clue that was related to three words, but I didn't consider the fact that the game master would stop the teammate's turn if they guess an opponent's word. This led to the teammate's turn ending prematurely. Next time, I will make sure to provide a clue that is related to only one or two words, so that the teammate's turn can continue even if they guess an opponent's word.} \\
\\ 
\texttt{Trial \#2: New plan: In my last attempt, I provided a clue that was related to three words, but I didn't consider the fact that the game master would stop the teammate's turn if they guess an opponent's word. This led to the teammate's turn ending prematurely. Next time, I will make sure to provide a clue that is related to only one or two words, so that the teammate's turn can continue even if they guess an opponent's word. However, I also made another mistake by providing a clue that was too obvious, which led to the teammate guessing an innocent word (pit) instead of a target word. Next time, I will try to provide a more subtle clue that is still related to the target words, but not too obvious.} \\
\\ 
\texttt{Trial \#3: New plan: In my last attempt, I provided a clue that was related to two words, but I didn't consider the fact that the game master would stop the teammate's turn if they guess an opponent's word, and I also didn't consider the fact that the teammate might guess an innocent word instead of a target word. This led to the teammate's turn ending prematurely. Next time, I will make sure to provide a clue that is related to only one word, so that the teammate's turn can continue even if they guess an opponent's word, and I will try to provide a more subtle clue that is still related to the target words, but not too obvious. I will also make sure to provide a clue that is semantically related to the target words, but not too obvious, and I will try to provide a clue that is related to a word that is not too easy to guess, but not too hard either.} \\
\\ 
\\ 
\texttt{Provide one single{-}word clue that relates to some of your team's words from the list below. You can choose to how many and to which words your clue relates to.} \\
\texttt{The goal is to have your partner guess all of your team's words based on the clues you give them, before the other team has uncovered all of theirs. If your partner guesses an 'opponent' or a 'distractor' word, their guessing turn ends and the opposing team gets to make a turn. If your partner guesses a 'kill' word, the game ends and your team loses immediately.} \\
\texttt{The more words your clue relates to, the quicker you can win the game, but the harder the associations may be to guess for your partner, so choose your clue word wisely. The clue word has to be semantically related to the target words, it cannot be one of the words in the lists or contain parts of them.} \\
\\ 
\texttt{Always give your single{-}word clue and your comma{-}separated list of related target words in the following format and make your answers as short as possible, never include any other text than is required in this form:} \\
\\ 
\texttt{CLUE: <WORD>} \\
\texttt{TARGETS: <TARGETS>} \\
\\ 
\texttt{Your team words are: angle, flush.} \\
\texttt{Your opponent's team words are: .} \\
\texttt{Distractor words are: cuff, boil, jam, fan, rocket.} \\
\texttt{The kill words are: beetle.} \\
            }
        }
    }
    & \\ \\

    \theutterance \stepcounter{utterance}  
    & \multicolumn{4}{p{0.6\linewidth}}{
        \cellcolor[rgb]{0.9,0.9,0.9}{
            \makecell[{{p{\linewidth}}}]{
                \texttt{\tiny{[P1$\rangle$GM]}}
                \texttt{CLUE: edge} \\
\texttt{TARGETS: angle, flush} \\
            }
        }
    }
    & & \\ \\

    \theutterance \stepcounter{utterance}  
    & & & \multicolumn{2}{p{0.3\linewidth}}{
        \cellcolor[rgb]{0.9,0.9,0.9}{
            \makecell[{{p{\linewidth}}}]{
                \texttt{\tiny{[GM$|$GM]}}
                \texttt{edge} \\
            }
        }
    }
    & & \\ \\

    \theutterance \stepcounter{utterance}  
    & & & \multicolumn{2}{p{0.3\linewidth}}{
        \cellcolor[rgb]{0.9,0.9,0.9}{
            \makecell[{{p{\linewidth}}}]{
                \texttt{\tiny{[GM$|$GM]}}
                ['angle', 'flush']
            }
        }
    }
    & & \\ \\

    \theutterance \stepcounter{utterance}  
    & & & \multicolumn{2}{p{0.3\linewidth}}{
        \cellcolor[rgb]{0.9,0.9,0.9}{
            \makecell[{{p{\linewidth}}}]{
                \texttt{\tiny{[GM$|$GM]}}
                {'word': 'angle', 'assignment': 'team'}
            }
        }
    }
    & & \\ \\

    \theutterance \stepcounter{utterance}  
    & & & \multicolumn{2}{p{0.3\linewidth}}{
        \cellcolor[rgb]{0.9,0.9,0.9}{
            \makecell[{{p{\linewidth}}}]{
                \texttt{\tiny{[GM$|$GM]}}
                {'word': 'flush', 'assignment': 'team'}
            }
        }
    }
    & & \\ \\

    \theutterance \stepcounter{utterance}  
    & & & \multicolumn{2}{p{0.3\linewidth}}{
        \cellcolor[rgb]{0.9,0.9,0.9}{
            \makecell[{{p{\linewidth}}}]{
                \texttt{\tiny{[GM$|$GM]}}
                \texttt{opponent has won} \\
            }
        }
    }
    & & \\ \\

    \theutterance \stepcounter{utterance}  
    & \multicolumn{4}{p{0.6\linewidth}}{
        \cellcolor[rgb]{0.95,0.95,0.95}{
            \makecell[{{p{\linewidth}}}]{
                \texttt{\tiny{[P2$\rightarrow$P1]}}
                \texttt{Post{-}episode Reflexion (not sent to GM):} \\
\texttt{New plan: In my last attempts, I provided clues that were too obvious, leading to the teammate guessing innocent words instead of target words. I also didn't consider the fact that the game master would stop the teammate's turn if they guess an opponent's word. Next time, I will try to provide a more subtle clue that is still related to the target words, but not too obvious. I will also make sure to provide a clue that is semantically related to the target words, but not too obvious, and I will try to provide a clue that is related to a word that is not too easy to guess, but not too hard either. Additionally, I will make sure to provide a clue that is related to only one word, so that the teammate's turn can continue even if they guess an opponent's word. To achieve this, I will carefully analyze the words in the list and look for a word that has a unique connection to one of the target words, without being too obvious or too hard to guess.} \\
            }
        }
    }
    & & \\ \\

\end{supertabular}
}

\end{document}
